% ID: SEL-D0014-BD-E08
% ESTADO: cotejada. Segunda pasada visual completa y borrador compilado revisados.
% Fuente: Selectividad/matematicas_aplicadas/matematicas_aplicadas_P_B.pdf
% SHA-256: 51601c4f5a6c1dfcd860f621a659f8b9c09f53ff4cd31a1fa1e827dbd8f06159
% Página: 2; bloque D; ejercicio original 8.

Se sabe que la longitud, en centímetros, de una especie de estrella de mar
en una determinada zona sigue una ley Normal con desviación típica 3. Para
estimar la longitud media de esa especie de estrella de mar, se extrae una
muestra de tamaño 36 y se obtiene el intervalo de confianza $(6.04,8)$ al
95\%. Se pide:
\begin{enumerate}[label=\alph*)]
  \item \textbf{(0.5 puntos)} Calcule la media muestral.
  \item \textbf{(0.5 puntos)} Calcule el error de estimación máximo cometido.
  \item \textbf{(1 punto)} Si aumentamos el tamaño muestral a 49,
  \textquestiondown{}qué efecto produce sobre el error máximo cometido?
  Calcule este error.
  \item \textbf{(0.5 puntos)} Si aumentamos el nivel de confianza,
  \textquestiondown{}qué efecto produce sobre el error de estimación máximo?
  Justifique la respuesta.
\end{enumerate}
